\scanpage{42}
ment'' evident, which here says \deleted{\unsure{that}}\inserted{that the datum}
of a gerbe \inserted{$\mathcal G_X$}
on $X$, amounts to that of \unsure{gerbes|gerbe} $\mathcal G_{X^*}$ on
$X^*$ and $\mathcal G_{Y,X}$ on $\mathcal V_{Y,X}$, and of an
\deleted{isomorphism}\inserted{equivalence} of \deleted{\unsure{gluing|glues|to glue}} gerbes
\[
(24)\qquad \mathcal G_{Y,X}\mid\mathcal V^*_{Y,X}\simeq
\mathcal G_{X^*}\mid\mathcal V^*_{Y,X}.
\]
Now on a $K(\pi,1)$, an $A$-gerbe
is interpreted canonically as an
\uline{extension} of $\pi$ by $A$. Thus the datum
\margin{\diagram{\begin{tikzcd}[column sep=small,row sep=small]
D\arrow[r]\arrow[d]&\Pi_1(Y)\\
\Pi_1(X^*)&
\end{tikzcd}}}
of a gerbe on $X\deleted{{}\unsure{\simeq} S^2}$\unsure{,} amounts to that
of \deleted{\unsure{an extension}} two extensions, one of
$\Pi_1(Y)$ by $A$, the other of $\Pi_1(X^*)$ by $A$
and of an isomorphism between the two extensions
of $D$ by $A$ which are deduced from them\unsure{,}
by inverse image under the homomorphisms (10).
(this description works as soon as the
``\unsure{stones|pieces|paving stones}'' $Y$, $X'$, $\mathcal V^*_{Y,X}$ are connected, and
their $\Pi_2$ is zero). In the case
% Pale lower-left foliation, another hand; not authorial page numbering.
\margin{41}
